\documentclass[11pt]{article}
\usepackage[margin=1in]{geometry}
\usepackage{amsmath,amssymb,amsthm,mathtools}
\usepackage{graphicx,xcolor,microtype,hyperref}
\hypersetup{colorlinks=true,linkcolor=blue!55!black,citecolor=blue!55!black,urlcolor=blue!55!black}
\newtheorem{theorem}{Theorem}
\newtheorem{lemma}{Lemma}
\newcommand{\T}{\mathbb T}
\newcommand{\D}{\mathbb D}
\newcommand{\C}{\mathbb C}
\newcommand{\one}{\mathbf 1}
\title{\vspace{-2em}Minimal Unitary Realizations Can Agree Locally\\
and Disagree on a Disconnected Component}
\author{Counterexample packet for arXiv:2509.08921}
\date{27 August 2026}

\begin{document}
\maketitle
\begin{center}
\fcolorbox{red!65!black}{red!3}{\parbox{0.91\textwidth}{
\textbf{Status: candidate full counterexample, likely valid.}
There are two minimal, analytically equivalent univariate realizations whose
operators are unitary with spectrum $\T$, but whose transfer functions differ
by exactly $-1/z$ on $\C\setminus\overline\D$. This supplies the example
requested in Remark~2.4 of the source paper.
}}
\end{center}

\section{Source and exact question}

The source is A.~Karoobi, R.~T.~W.~Martin, and M.~Tornes,
\emph{Operator realizations about a matrix-centre}, arXiv:2509.08921
\cite{KMT}. For a bounded operator $A$ on a Hilbert space and vectors $b,c$,
the scalar transfer function is
\[
                  f_A(z)=b^*(I-zA)^{-1}c.
\]
The realization is \emph{minimal} when $c$ is $A$-cyclic and $b$ is
$A^*$-cyclic. Two realizations are analytically equivalent at $0$ exactly when
their moments $b^*A^n c$ agree for every $n\geq0$; equivalently, their
matrix-level transfer power series agree near $0$.

Remark~2.4 on source PDF page 8 asks whether there exist two minimal,
analytically equivalent realizations with disconnected invertibility domains
which define different analytic functions on a component disconnected from
$0$.

\begin{figure}[ht]
\centering
\includegraphics[width=0.98\textwidth]{figures/open_problem_crop.png}
\caption{Source PDF page 8: Example~2.3 and the exact question in
Remark~2.4.}
\end{figure}

\section{Proof intuition}

The source's bilateral-shift example fails minimality because a point mass at
one Fourier mode sees only half of the bilateral orbit. We replace that point
mass by Lebesgue measure restricted to a proper measurable part of the unit
circle. On any proper piece whose complement has positive measure, Hardy
boundary uniqueness makes both analytic and anti-analytic monomials dense;
this gives the two cyclicity conditions and hence minimality.

Split the circle into two pieces $E$ and $F$, each occurring in every arc.
The first realization uses $L^2(F)$ and output vector $\bar\zeta$, while the
second uses $L^2(E)$ and output vector $-\bar\zeta$. Their nonnegative moments
add up to the positive Fourier coefficients of Lebesgue measure, hence cancel.
On the exterior component the missing Laurent coefficient is the constant
Fourier coefficient instead, and it survives as $-1/z$.

\section{Two elementary lemmas}

Let $m$ denote normalized Lebesgue measure on $\T$.

\begin{lemma}[Everywhere-thick measurable partition]\label{lem:partition}
There is a measurable partition $\T=E\mathbin{\dot\cup}F$ such that
\[
             m(E\cap I)>0\quad\text{and}\quad m(F\cap I)>0
\]
for every nonempty open arc $I\subset\T$.
\end{lemma}

\begin{proof}
Enumerate a countable basis of open arcs as $(I_k)_{k\geq1}$. Inductively,
inside $I_k$ and away from the finitely many compact sets already selected,
choose two disjoint compact nowhere-dense sets $K_k,L_k$ of positive measure.
This is possible by placing two small fat Cantor sets in an open subarc of the
remaining set. Put $E=\bigcup_{k\geq1}K_k$ and $F=\T\setminus E$. Every future
$K_j$ is chosen disjoint from every earlier $L_k$, so $L_k\subset F$. Every
open arc contains some basis arc $I_k$ and therefore contains positive-measure
subsets $K_k\subset E$ and $L_k\subset F$.
\end{proof}

\begin{lemma}[One-sided polynomial density]\label{lem:density}
If $G\subset\T$ is measurable and $m(\T\setminus G)>0$, then both
\[
 \operatorname{span}\{\zeta^n:n\geq0\}
 \quad\text{and}\quad
 \operatorname{span}\{\bar\zeta^n:n\geq0\}
\]
are dense in $L^2(G,m)$.
\end{lemma}

\begin{proof}
Suppose $h\in L^2(G)$ is orthogonal to every $\zeta^n$, $n\geq0$. Extend
$\one_G\bar h$ by zero to $\T$. Its Fourier coefficients of index at most zero
vanish, so it is the boundary function of an element of $zH^2$. It vanishes on
$\T\setminus G$, a set of positive measure. Hardy boundary uniqueness forces
it to vanish identically; see, for example, \cite[Chapter~II]{Duren}. Hence
$h=0$. This proves density of analytic polynomials. Complex conjugation gives
the anti-analytic assertion.
\end{proof}

\section{The counterexample}

Choose $E,F$ as in Lemma~\ref{lem:partition}. For $G\in\{E,F\}$ let
\[
             \mathcal H_G=L^2(G,m),\qquad
             (U_Gh)(\zeta)=\zeta h(\zeta).
\]
Define two realizations by
\begin{equation}\label{eq:triples}
 (A,b,c)=(U_F,\bar\zeta,\one_F),
 \qquad
 (A',b',c')=(U_E,-\bar\zeta,\one_E).
\end{equation}

\begin{theorem}\label{thm:main}
The two realizations in \eqref{eq:triples} are minimal and analytically
equivalent at $0$. Both scalar invertibility domains are
\[
                \C\setminus\T=\D\mathbin{\dot\cup}
                (\C\setminus\overline\D).
\]
If $f_F$ and $f_E$ are their scalar transfer functions, then
\[
 f_F(z)=f_E(z)\quad(|z|<1),
 \qquad
 f_F(z)-f_E(z)=-\frac1z\quad(|z|>1).
\]
Thus they give the example requested in Remark~2.4 of \cite{KMT}.
\end{theorem}

\begin{proof}
Each $U_G$ is unitary. Since every open arc has positive $G$-measure, the
essential range of the coordinate function on $G$ is all of $\T$. Thus
$\sigma(U_E)=\sigma(U_F)=\T$, and $I-zU_G$ is invertible exactly when
$|z|\neq1$.

By Lemma~\ref{lem:density}, $\one_G$ is $U_G$-cyclic. The adjoint is
$U_G^*=M_{\bar\zeta}$, and
\[
        (U_F^*)^n\bar\zeta=\bar\zeta^{n+1},\qquad
        (U_E^*)^n(-\bar\zeta)=-\bar\zeta^{n+1}.
\]
The shifted anti-analytic monomials are dense because multiplication by
$\bar\zeta$ is unitary and the unshifted anti-analytic monomials are dense by
Lemma~\ref{lem:density}. Hence both realizations are minimal.

For every $n\geq0$, their Markov parameters are
\begin{align*}
 b^*A^n c&=\int_F \zeta^{n+1}\,dm(\zeta),\\
 b'^*A'^n c'&=-\int_E \zeta^{n+1}\,dm(\zeta).
\end{align*}
Because $E\mathbin{\dot\cup}F=\T$ modulo null sets,
\[
 b^*A^n c-b'^*A'^n c'
   =\int_{\T}\zeta^{n+1}\,dm(\zeta)=0.
\]
The scalar transfer series therefore agree on $\D$. More strongly, at every
matrix level and for every matrix $Z$ with $\lVert Z\rVert<1$, both quantized
transfer functions are
\[
                 \sum_{n\geq0} Z^n\,(b^*A^n c),
\]
so the realizations are analytically equivalent at $0$ in the precise sense
of the source paper.

Finally, for any $|z|\neq1$,
\begin{align*}
 f_F(z)-f_E(z)
 &=\int_F\frac{\zeta}{{1-z\zeta}}\,dm(\zeta)
   +\int_E\frac{\zeta}{1-z\zeta}\,dm(\zeta)\\
 &=\int_{\T}\frac{\zeta}{1-z\zeta}\,dm(\zeta).
\end{align*}
For $|z|<1$, the geometric series has only strictly positive Fourier modes,
so the last integral is zero. For $|z|>1$,
\[
 \frac{\zeta}{1-z\zeta}
   =-\frac1z\sum_{k=0}^{\infty}z^{-k}\zeta^{-k},
\]
uniformly on $\T$, and integration leaves only the $k=0$ term. The difference
is therefore $-1/z$, as claimed.
\end{proof}

\section{Verification, novelty, and scope}

\paragraph{Verification.}
The accompanying verification report audits separately the measurable
partition, the two cyclicity statements, the spectrum calculation, equality
at every matrix level, and the exterior Laurent coefficient. No numerical
computation is used; all identities are exact.

\paragraph{Bounded novelty search.}
The four run indexes were searched for arXiv:2509.08921, the exact title,
``minimal and analytically equivalent realizations,'' disconnected
invertibility domains, unitary realizations, and Markov parameters.
Exact-phrase and close arXiv/web searches through 27 August 2026 returned the
source and general realization literature, but no later paper or preprint
answering Remark~2.4. Novelty is plausible, not certified.

\paragraph{Scope and review.}
This is a full affirmative answer to the existence question in Remark~2.4:
both realizations are minimal, both domains have the required disconnected
component, and their global values differ there. It does not address the
separate compact matrix-centre realization question in Section~9 of the source.
Human review should focus on Lemma~\ref{lem:density} and the sign convention in
the two moment formulas.

\begingroup
\sloppy
\raggedright
\begin{thebibliography}{9}
\bibitem{KMT}
A.~Karoobi, R.~T.~W.~Martin, and M.~Tornes,
\emph{Operator realizations about a matrix-centre},
\href{https://arxiv.org/abs/2509.08921}{arXiv:2509.08921}, 2025.

\bibitem{Duren}
P.~L.~Duren,
\emph{Theory of $H^p$ Spaces},
Pure and Applied Mathematics 38, Academic Press, 1970;
Dover reprint, 2000.
\end{thebibliography}
\endgroup
\end{document}
